\subsubsection{对椭圆 PDE 的应用}\label{subsec:implicit-function-elliptic-applications}

希望将前一小节的结果应用于非线性椭圆 PDE. 为此, 必须将这些问题写成 Banach 空间中的抽象非线性方程. 以下引理对此至关重要.

\begin{Lemma}
\label{lem:sobolev-pointwise-composition-differentiability}
设 $\Omega$ 为 $\mathbb{R}^n$ 中具有光滑边界的有界区域. 再设 $f:\mathbb{R}^k\to\mathbb{R}$ 属于 $C^{m+1}$ 类, 其中 $m>n/2$. 则将 $f$ 逐点作用, 会诱导从 $(H^m(\Omega))^k$ 到 $H^m(\Omega)$ 的连续可微映射.
\end{Lemma}

\begingroup
\renewcommand{\endProof}{\par\nobreak\hfill$\blacksquare$\par\vspace{0.1cm}}
\begin{Proof}
由 Sobolev 嵌入定理, 有 $H^m(\Omega)\subset C(\overline\Omega)$. 因而对每个 $\bm{u}\in(H^m(\Omega))^k$, 都有 $f(\bm{u})\in C(\overline\Omega)\subset L^2(\Omega)$. 此外, 若 $\bm{u}$ 属于 $(C^m(\overline\Omega))^k$, 则可由链式法则求出 $f(\bm{u})$ 的导数; 一般情形则可利用光滑函数逼近. 例如, 设 $\bm{u}=\lim_{n\to\infty}\bm{u}_n$ 在 $(H^m(\Omega))^k$ 的拓扑下成立, 其中 $\bm{u}_n$ 光滑. 得到
\begin{equation}\label{eq:sobolev-composition-distributional-chain-rule}
\begin{aligned}
\left(f(\bm{u}),\frac{\partial\phi}{\partial x_i}\right)
&=\lim_{n\to\infty}\left(f(\bm{u}_n),\frac{\partial\phi}{\partial x_i}\right)\\
&=-\lim_{n\to\infty}\left(\nabla f(\bm{u}_n)\frac{\partial\bm{u}_n}{\partial x_i},\phi\right)\\
&=-\left(\nabla f(\bm{u})\frac{\partial\bm{u}}{\partial x_i},\phi\right),
\end{aligned}
\end{equation}
即链式法则适用于 $f(\bm{u})$ 的求导. 类似地可以处理高阶导数. 注意 $f(\bm{u})$ 的所有导数都具有乘积形式, 其中含有 $f$ 的一个导数和 $\bm{u}$ 的若干导数. 第一因子属于 $C(\overline\Omega)$, 而 $\bm{u}$ 的任意 $l$ 阶导数属于 $H^{m-l}(\Omega)$; 若 $m-l<n/2$, 后者嵌入 $L^{2n/(n-2(m-l))}(\Omega)$. 利用这一事实和 Hölder 不等式, 可以证明 $f(\bm{u})$ 的所有不超过 $m$ 阶的导数都属于 $L^2(\Omega)$; 此外, 由这一论证显然可知 $f$ 实际上从 $(H^m(\Omega))^k$ 到 $H^m(\Omega)$ 连续. 利用关系
\begin{equation}\label{eq:sobolev-composition-integrated-difference-formula}
f(\bm{u})-f(\bm{v})=\int_0^1\nabla f(\bm{v}+\theta(\bm{u}-\bm{v}))\cdot(\bm{u}-\bm{v})\,\mathrm{d}\theta.
\end{equation}
可沿类似思路得到 $f$ 的可微性. 证明中更多的细节留给读者补充 (Problem~\ref{ex:sobolev-composition-lemma-details}).
\end{Proof}
\endgroup
